\begin{abstract}
Large unsupervised language models (LMs) possess extensive general knowledge and demonstrate some reasoning abilities, but exerting fine-grained control over their outputs remains a challenge, primarily because their training lacks supervision. Current approaches to enhance steerability typically involve collecting human judgements comparing model outputs and subsequently fine-tuning the LM to reflect these preferences, frequently employing reinforcement learning from human feedback (RLHF). Nonetheless, RLHF entails a multi-stage and sometimes unstable process: it first constructs a reward model that encapsulates human preferences, then applies reinforcement learning to adapt the LM, aiming to maximize the learned reward function while minimizing divergence from the pretrained model. In this work, we propose a novel formulation of the reward model within the RLHF framework that admits a closed-form solution for the associated optimal policy, allowing the standard RLHF objective to be addressed with a straightforward classification loss. We term this approach \textit{Direct Preference Optimization} (DPO), which offers improved stability, efficiency, and ease of use by removing the requirement to sample from the LM during fine-tuning or to conduct extensive hyperparameter searches. Experimental results demonstrate that DPO is at least as effective as, and often outperforms, existing techniques in aligning LMs with human preferences. In particular, DPO fine-tuning provides better control over generation sentiment than PPO-based RLHF, and achieves equivalent or superior results in tasks like summarization and single-turn dialogue, all while considerably simplifying both implementation and training. 
\end{abstract}

\section{Introduction}
Large unsupervised language models (LMs), trained on massive corpora, have been found to exhibit unexpected and impressive capabilities~\citep{chowdhery2022palm, brown2020language, touvron2023llama,bubeck2023sparks}. Despite these achievements, the training data for such models originates from human-authored sources that reflect a broad spectrum of intentions, values, and expertise. Consequently, some of the behaviors and competencies encoded in these models may be undesirable for emulation. For instance, it is beneficial for an AI programming assistant to \textit{recognize} prevalent programming errors to aid in their correction; however, when the assistant produces code, we prefer its outputs to be influenced by the minority instances of superior programming found within the data, rather than reproducing common mistakes. Likewise, while a language model should be \textit{cognizant} of widespread misconceptions (such as those endorsed by half the population), we would not wish it to replicate these inaccuracies in half of the relevant responses it generates. Thus, a critical challenge lies in isolating and amplifying a model’s \emph{intended responses and behaviors} from its expansive reservoir of \textit{knowledge and competencies}, in pursuit of creating AI systems that are safe, robust, and manageable \citep{ouyang2022training}. Standard strategies for achieving this objective typically direct language models using reinforcement learning (RL) aligned with human values, yet we demonstrate that the RL-based goals traditionally employed can, in fact, be realized precisely through a straightforward binary cross-entropy objective. This approach substantially streamlines the process of learning from preferences.

Broadly, current methodologies inculcate language models with the appropriate behaviors by leveraging carefully collected datasets that reflect human value judgments on what constitutes helpful and safe outputs. This process of preference learning is implemented following an initial large-scale unsupervised pre-training phase using diverse textual data. The most basic technique for aligning model outputs with human preferences involves supervised fine-tuning with high-quality human-generated exemplars. Nevertheless, reinforcement learning from human (or AI) feedback (RLHF/RLAIF; \citep{christiano2017deep,bai2022constitutional}) has proven to be the most effective class of alignment techniques. RLHF consists of training a reward model on preference data, then applying RL to update the language model so that it produces responses that receive higher rewards, all while constraining changes to stay close to the pre-trained model. Although RLHF yields models capable of sophisticated conversation and coding, its adoption introduces a much higher degree of procedural complexity relative to supervised approaches, including the need for training several language models and on-policy sampling throughout training—factors that collectively result in substantially increased computational demands.

This work demonstrates a method to optimize language models according to human preferences without necessitating reward modeling or the use of reinforcement learning techniques. We introduce \textit{{Direct Preference Optimization} (DPO)}, an algorithm that implicitly maximizes the same objective as traditional RLHF frameworks—namely, reward maximization regularized by a KL-divergence term—while being notably simpler to both implement and train. At a high level, DPO works by amplifying the log probability ratio between preferred and non-preferred outputs, but incorporates an adaptive, instance-specific weighting scheme that mitigates the model collapse often observed with naive probability ratio approaches. DPO, much like previous approaches, relies on a theoretical framework for modeling preferences—such as the Bradley-Terry model (\cite{bradley1952rankanalysis})—to gauge the correspondence between a candidate reward function and observed human preference data. Distinct from prior methods that utilize this preference model to construct a loss for reward model training—before then optimizing a policy against the resulting learned rewards—DPO reparameterizes the loss to be a direct function of the policy itself. Thus, equipped with a dataset of human feedback indicating preference between candidate outputs, DPO enables direct policy training via a straightforward binary cross-entropy loss, ultimately yielding the policy that is optimal for a latent reward function consistent with the empirical preference data.

The principal contribution presented here is Direct Preference Optimization (DPO), a lightweight alternative to RL-based approaches for training language models on preference data. Empirical results indicate that DPO matches or surpasses the performance of prevalent methods—including PPO-driven RLHF—across a range of tasks such as sentiment adjustment, summarization, and open-ended dialogue, utilizing language models with parameter counts up to 6B.

\section{Related Work}

Self-supervised language models, when scaled up, have demonstrated the ability to perform certain tasks in a zero-shot setting \citep{radford2019language} or with minimal in-context examples \citep{gpt3,megatron,chowdhery2022palm}. Nonetheless, further improvements in downstream task performance and alignment with user intent are achieved through fine-tuning on datasets composed of instructions and completions authored by humans \citep{mishra-etal-2022-cross,sanh2022multitask,chung2022scaling,thoppilan2022lamda}. Through this process—known as `instruction-tuning'—LLMs exhibit enhanced generalization to instructions not seen during fine-tuning, which contributes to their broader practical utility \citep{chung2022scaling}. While instruction tuning has shown significant promise, it is often more feasible to obtain \textit{comparative} human assessments of output quality than detailed expert demonstrations. In response, later research has involved fine-tuning LLMs on datasets reflecting human preference judgments, which has resulted in notable advances in tasks such as translation \citep{kreutzer-etal-2018-reliability}, summarization \citep{stiennon2022learning,ziegler2020finetuning}, story generation \citep{ziegler2020finetuning}, and following instructions \citep{ouyang2022training,ramamurthy2023is}. In these approaches, a reward function—typically represented by a neural network—is trained to be consistent with the dataset of human preferences under a model such as the Bradley-Terry model \citep{bradley1952rankanalysis}, after which the language model undergoes further training to maximize this reward using reinforcement learning algorithms, most commonly REINFORCE \citep{williams1992reinforce}, proximal policy optimization (PPO; \cite{schulman2017proximal}), or related variants \citep{ramamurthy2023is}. Related strands of work use LLMs, already instruction-tuned with human feedback, to autonomously generate new synthetic preference data, targeting characteristics like safety or harmlessness \citep{bai2022constitutional}, with only high-level human supervision via rubrics for annotation. Such techniques reflect a fusion of research trajectories: one focused on leveraging reinforcement learning to train language models towards various goals~\citep{Ranzato2015SequenceLT,paulus2018a,wu2018learning}, and another on developing general strategies for learning from human preferences \citep{christiano2017deep,kupcsik2018learning}. Although leveraging comparative human preferences is intuitively attractive, applying reinforcement learning for fine-tuning large language models introduces substantial practical difficulties; this work proposes a theoretically-sound methodology for optimizing relative preference data without the reliance on RL.

Beyond the realm of language modeling, the problem of deriving policies from preference information has been explored within both bandit and reinforcement learning frameworks, with numerous methodologies having been introduced. In particular, contextual bandit problems that rely on action preferences or rankings—rather than explicit reward feedback—are referred to as contextual dueling bandits (CDB; \cite{yue2012karmed,dudik2015contextual}). When explicit reward signals are unavailable, theoretical work in CDBs adopts the \textit{von Neumann winner} as a substitute for the conventional optimal policy; specifically, such a policy guarantees an expected win rate of at least 50\% against any alternative policy \citep{dudik2015contextual}. A key distinction is that, whereas preference annotations in CDBs are typically provided in an online manner, human preference learning often relies on a static, offline dataset containing preference-labeled action pairs \citep{yan2022human}. A related line of work, \textit{preference-based RL} (PbRL), formulates the task using binary preferences informed by an unobserved ‘scoring’ function in lieu of direct reward signals \citep{BusaFekete2014,ruiz2023dueling}. Numerous PbRL algorithms exist—some capable of utilizing off-policy preference datasets—but these generally involve an explicit intermediate step of inferring the hidden scoring function (i.e., constructing a reward model), followed by its subsequent optimization \citep{jain2013learning,BusaFekete2014,christiano2017deep,sadigh2017active,kupcsik2018learning}. Departing from this two-step paradigm, we propose a single-stage approach in which policy parameters are optimized directly to adhere to observed preferences.

\section{Preliminaries}\label{section:prelims}

Here, we summarize the RLHF pipeline as described in \citeauthor{ziegler2020finetuning}, and subsequently in works such as \citep{stiennon2022learning, bai2022training, ouyang2022training}. The standard process is composed of three sequential stages: (1) supervised fine-tuning (SFT), (2) collection of preference data and learning a reward model, and (3) optimization of the policy using reinforcement learning.

\textbf{SFT}: The RLHF workflow typically commences by adapting a large pre-trained language model to specific tasks (e.g., dialogue, summarization) through supervised learning on curated high-quality datasets, resulting in a fine-tuned policy $\pi^\text{SFT}$.

\textbf{Reward Modelling Phase}: The subsequent stage involves generating candidate response pairs by conditioning $\pi^\text{SFT}$ on input prompts $x$, which yields answer pairs $(y_1, y_2)\sim \pi^\text{SFT}(y \mid x)$. Human annotators are then tasked with expressing a preference over each pair, recording which completion is favored, denoted by $y_w\succ y_l \mid x$, where $y_w$ and $y_l$ signify the selected and non-selected completions, respectively. These preference annotations are presumed to reflect the values of an underlying, unobservable reward function $r^*(y, x)$. Several statistical models for preference data are available; among them, the Bradley-Terry (BT) model \cite{bradley1952rankanalysis} is prevalent, although the more flexible Plackett-Luce models \citep{plackett1975analysis, luce2012individual} can also be incorporated—particularly when multiple candidate completions are ranked. In the BT framework, the likelihood of a human expressing a given preference $p^*$ can be characterized as:

\begin{equation}\label{eq:bradley-terry}
    p^*(y_1\succ y_2 \mid x)=\frac{\exp\left(r^*(x, y_1)\right)}{\exp\left(r^*(x, y_1)\right) + \exp\left(r^*(x, y_2)\right)}.
\end{equation}

Given a fixed dataset of pairwise preference samples $\mathcal{D}=\bigl\{x^{(i)}, y_w^{(i)}, y_l^{(i)}\bigr\}_{i=1}^N$ drawn according to $p^*$, we can introduce a parameterized reward function $r_{\phi}(x, y)$ and fit its parameters using maximum likelihood estimation. Viewing this setup as a binary classification task, the associated negative log-likelihood loss takes the following form:

\begin{equation}\label{eq:reward_model}
    \mathcal{L}_R(r_{\phi}, \mathcal{D}) = -\mathbb{E}_{(x, y_w, y_l)\sim \mathcal{D}}\bigl[\log \sigma(r_{\phi}(x, y_w)- r_{\phi}(x, y_l))\bigr]
\end{equation}

Here, $\sigma$ denotes the sigmoid activation. When dealing with LMs, the reward network $r_{\phi}(x, y)$ is typically initialized from the supervised fine-tuning model $\pi^\text{SFT}(y \mid x)$, augmented with a linear head after the final transformer layer to yield a scalar reward prediction \cite{ziegler2020finetuning}. To reduce reward variance, it is common to apply normalization so that $\mathbb{E}_{x,y\sim \mathcal{D}}\left[r_\phi(x, y)\right] = 0$ holds for all $x$.

\textbf{RL Fine-Tuning Phase}: In the RL stage, the trained reward model provides evaluative signals to guide language model updates. Consistent with earlier studies~\citep{jaques2017sequence, jaques2020human}, the optimization procedure is expressed as

\begin{equation}\label{eq:RL}
\max_{\pi_{\theta}}  \mathbb{E}_{x\sim \mathcal{D}, y\sim \pi_{\theta}(y \mid x)}\bigl[r_{\phi}(x, y)\bigr] - \beta\mathbb{D}_{\textrm{KL}}\bigl[\pi_{\theta}(y\mid x)\mid \mid \pi_\text{ref}(y\mid x)\bigr],
\end{equation}

where $\beta$ acts as a regularization factor constraining deviation from the reference policy $\pi_\text{ref}$, typically taken as the original supervised model $\pi^\text{SFT}$. In practice, $\pi_\theta$ is also initialized from $\pi^\text{SFT}$. The KL regularization is crucial, as it ensures the policy remains in the regime where the reward model is reliable, preserves output diversity, and mitigates against collapse to a limited set of answers. Because text generation is discrete and thus non-differentiable, this objective is generally tackled using reinforcement learning algorithms. A common technique \citep{ziegler2020finetuning, stiennon2022learning, bai2022training, ouyang2022training} is to formulate the reward as ${r(x, y) = r_{\phi}(x, y) -\beta (\log \pi_{\theta}(y\mid x) - \log \pi_\text{ref}(y\mid x))}$, and maximize it using PPO \cite{schulman2017proximal}.

\section{Direct Preference Optimization}\label{sec:DPO}

Given the complexity and scale of applying RL to fine-tune large language models, we seek an alternative, more straightforward approach for policy optimization that utilizes preference data directly. Distinct from established RLHF pipelines, which require a separate reward model and subsequent RL optimization, our methodology exploits a specific reward model parameterization that admits a closed-form solution for the corresponding optimal policy—eliminating the need for an explicit RL training cycle. In what follows, we detail how this key idea enables the construction of an analytic mapping from rewards to optimal policies. This, in turn, allows us to recast objectives over reward functions into equivalent losses over the policy space. Through this change-of-variables strategy, our method dispenses with training a separate reward model but still operates compatibly with well-established preference models such as the Bradley-Terry framework. Functionally, the policy network now embodies both the gener

\textbf{Derivation of the DPO Objective.} We begin with the reinforcement learning objective presented in Eq.~\ref{eq:RL}, adopting an arbitrary reward function $r$. Building on earlier research~\citep{peters2007reinforcement, peng2019advantage, korbak2022reinforcement, go2023aligning}, it is readily demonstrated that the optimal policy for the reward maximization objective with a KL constraint, as stated in Eq.~\ref{eq:RL}, is given by:

\begin{equation}\label{eq:op_policy}
    \pi_r(y\mid x) = \frac{1}{Z(x)}\pi_\text{ref}(y\mid x)\exp\left(\frac{1}{\beta}r(x, y)\right),
\end{equation}

where $Z(x) =\sum_{y}\pi_\text{ref}(y\mid x)\exp\left(\frac{1}{\beta}r(x, y)\right)$ denotes the partition function. A detailed derivation is provided in Appendix~\ref{app:derivation1}. Direct computation of $Z(x)$ can be computationally demanding, even if we employ the MLE estimator $r_{\phi}$ for the true reward function $r^*$, as also pointed out in~\citep{korbak2022reinforcement, go2023aligning}. This complexity poses challenges for practical implementation of the above form. Nonetheless, Eq.~\ref{eq:op_policy} can be manipulated algebraically to write the reward in terms of its optimal policy $\pi_r$, the reference policy $\pi_\text{ref}$, and the yet-unknown partition function $Z(\cdot)$. Taking the logarithm of both sides of Eq.~\ref{eq:op_policy} and rearranging yields:

\begin{equation}\label{eq:main_eq}
    r(x,y) =\beta \log \frac{\pi_r(y\mid x)}{\pi_\text{ref}(y\mid x)} + \beta \log Z(x).
\end{equation}

This transformation can be equivalently applied to the ground-truth reward $r^*$ as well as the corresponding optimal policy $\pi^*$. Importantly, the Bradley-Terry model depends solely on the difference in rewards assigned to two outputs: specifically, ${p^*(y_1 \succ y_2 \mid x) = \sigma(r^*(x, y_1) - r^*(x, y_2))}$. Substituting the parameterization from Eq.~\ref{eq:main_eq} for $r^*(x,y)$ into the preference model given by Eq.~\ref{eq:bradley-terry}, the terms involving the partition function $Z(x)$ cancel out. This results in a human preference likelihood that depends only on the optimal policy $\pi^*$ and the reference policy $\pi_\text{ref}$. Thus, under the Bradley-Terry model, the RLHF-optimal policy $\pi^*$ fulfills the preference-based criterion:

\begin{equation}\label{eq:objective}
    p^*(y_1\succ y_2 \mid x)=\frac{1}{1 + \exp\left(\beta \log \frac{\pi^*(y_2\mid x)}{\pi_\text{ref}(y_2\mid x)} - \beta \log \frac{\pi^*(y_1\mid x)}{\pi_\text{ref}(y_1\mid x)}\right)}
\end{equation}

A full derivation appears in Appendix~\ref{app:derivation2}. While this formulation leverages the Bradley-Terry framework, analogous derivations may be obtained for more general models such as Plackett-Luce~\citep{plackett1975analysis, luce2012individual}, as detailed in Appendix~\ref{app:plackett_luce_models}.

Expressing the human preference probabilities directly in terms of the policy, rather than a reward model, allows us to construct a maximum likelihood estimation procedure for a parametrized policy $\pi_\theta$. In analogy to the reward modeling formulation (see Eq.~\ref{eq:reward_model}), we propose the following objective for direct policy optimization:

\begin{equation}\label{eq:optimum_model}
    \mathcal{L}_\text{DPO}(\pi_{\theta}; \pi_\text{ref}) = -\mathbb{E}_{(x, y_w, y_l)\sim \mathcal{D}}\left[\log \sigma \left(\beta \log \frac{\pi_{\theta}(y_w\mid x)}{\pi_\text

By employing this approach, we construct an implicit reward through a different parameterization, such that its corresponding optimal policy is directly represented by $\pi_\theta$. Additionally, since the approach aligns with training a reparametrized Bradley-Terry model, it inherits certain desirable theoretical guarantees, including consistency under appropriate assumptions about the distribution of the preference data \cite{bong2022generalized}. Further theoretical implications of DPO and its connections to related methods are elaborated in Section~\ref{sec:theory}.

\textbf{Mechanistic function of the DPO update.} To gain deeper intuition for how DPO operates, we examine the gradient of its objective, $\mathcal{L}_\text{DPO}$. The gradient with respect to $\theta$ is given by:

\begin{multline*}\label{eq:gradient}
    \nabla_\theta \mathcal{L}_\text{DPO}(\pi_\theta;\pi_\text{ref}) = \\ -\beta\mathbb{E}_{(x, y_w, y_l) \sim \mathcal{D}} \bigg[\underbrace{\sigma(\hat{r}_\theta(x, y_l) - \hat{r}_\theta (x, y_w))}_\text{larger impact when reward assignment is inaccurate}\bigg[\underbrace{\nabla_\theta\log \pi(y_w \mid x)}_\text{push up likelihood of $y_w$} - \underbrace{\nabla_\theta\log\pi(y_l \mid x)}_\text{push down likelihood of $y_l$}\bigg]\bigg],
\end{multline*}

where the implicit reward assigned by the model, $\hat{r}_\theta(x, y)$, is defined as $\beta \log \frac{\pi_\theta(y \mid x)}{\pi_\text{ref}(y \mid x)}$ (see further details in Section~\ref{sec:theory}). In essence, the gradient of $\mathcal{L}_\text{DPO}$ functions to elevate the probability of favorable completions $y_w$ while reducing the probability of unfavorable ones $y_l$. Notably, each training example is reweighted according to the extent to which the implicit reward $\hat{r}_\theta$ overvalues the dispreferred response, modulated by $\beta$—that is, by the magnitude of error in the implicit reward's ordering, which incorporates the effect of the KL regularization. Empirical findings indicate that this reweighting is crucial; omitting the coefficient—thus applying a uniform weighting—tends to lead the language model to collapse, as seen in Appendix Table~\ref{tab:unlikelihood_generations}.

\textbf{DPO Overview.} The standard DPO process involves the following steps: 1) For each prompt $x$, generate completions $y_1, y_2$ by sampling from $\pi_\text{ref}(\cdot \mid x)$ and annotate them using human judgments to form an offline preference dataset $\mathcal{D} = \{x^{(i)}, y_w^{(i)}, y_l^{(i)}\}_{i=1}^N$; 2) update the language model $\pi_\theta$ by minimizing $\mathcal{L}_\text{DPO}$ given $\pi_\text{ref}$, $\mathcal{D}$, and the specified $\beta$. In application, leveraging existing public preference datasets is desirable to avoid manual generation and labeling. Since these datasets are typically generated using $\pi^\text{SFT}$, we set $\pi_\text{ref} = \pi^\text{SFT}$ when it is accessible. In scenarios where $\pi^\text{SFT}$ is not provided, $\pi_\text{ref}$ is initialized by maximizing the likelihood over preferred completions ${(x, y_w)}$, specifically, by setting ${\pi_\text{ref} = \argmax_{\pi}\mathbb{E}_{x, y_w \sim \mathcal{D}}[\log \pi(y_w \mid x)]}$. This method helps address the mismatch between the actual reference distribution, which is unattainable, and the approximation $\pi_\text{ref}$ adopted in DPO. Additional implementation specifics and choices of hyperparameters are documented in Appendix~\ref{app:implementation}.

\section{Theoretical Analysis of DPO}
Here, we provide a deeper theoretical understanding of the DPO approach, establishing its foundations and discussing how DPO addresses certain limitations observed in actor-critic methods employed for RLHF—such as PPO~\cite{schulman2017proximal}.

\label{sec:theory}

\subsection{Your Language Model Is Secretly a Reward Model}
DPO facilitates learning by avoiding explicit reward function estimation and separate reinforcement learning steps; instead, it directly uses a single maximum likelihood objective. The objective from Eq. \ref{eq:main_eq} corresponds to the Bradley-Terry framework, wherein rewards are parameterized as $r^*(x, y) = \beta \log\frac{\pi^*_\theta(y \mid x)}{\pi_\text{ref}(y \mid x)}$. The parameterized model $\pi_{\theta}$ is trained in an analogous fashion to the reward model objective in Eq. \ref{eq:reward_model}, utilizing a change of variables. This section develops the theoretical perspective underlying this reparameterization, demonstrates that it does not restrict the class of potential reward models, and proves it permits full recovery of the optimal policy. We initiate this analysis by establishing an equivalence relation among reward functions.

\begin{definition}
Two reward functions $r(x, y)$ and $r'(x, y)$ are called equivalent if and only if there exists a function $f$ such that ${r(x, y)-r'(x, y) = f(x)}$.
\end{definition}

This definition immediately implies an equivalence relation, segmenting the entire collection of reward functions into distinct equivalence classes. The following lemmas capture key properties:

\begin{lemma}\label{lemma:same_prefrence}
Within the Plackett-Luce framework, including the Bradley-Terry model as a special case, any pair of reward functions from the same equivalence class induce identical preference distributions.
\end{lemma}

\begin{lemma}\label{lemma:same_policy}
    Any two reward functions belonging to a shared equivalence class yield the same optimal policy in the context of the constrained RL problem.
\end{lemma}

The arguments supporting these lemmas are elementary and thus relegated to Appendix \ref{app:lemma1}. The first lemma highlights a classic under-specification challenge inherent to Plackett-Luce models \cite{plackett1975analysis}. This lack of identifiability necessitates imposing supplementary constraints in practice to ensure meaningful MLE results from Eq.~\ref{eq:reward_model} \cite{bong2022generalized}. The second lemma indicates that all rewards in an equivalence class produce the same optimal policy; thus, when addressing the ultimate objective, it suffices to recover any representative from the correct class. We formalize this insight in the following theorem, whose proof can be found in Appendix~\ref{app:thm1}:

\begin{theorem}\label{thm:main}
    Provided mild regularity conditions, every equivalence class of reward functions compatible with the Plackett-Luce (and particularly the Bradley-Terry) family can be described by the reparameterization ${r(x, y) = \beta \log \frac{\pi(y\mid x)}{\pi_\text{ref}(y\mid x)}}$ for some choice of model $\pi(y\mid x)$ and a designated reference policy $\pi_\text{ref}(y \mid x)$.
\end{theorem}

\begin{sproof}
    Let $r(x, y)$ be a given reward function, inducing an optimal policy $\pi_r(y \mid x)$ as defined by Eq.~\ref{eq:op_policy}. We aim to show that for any such $r$, an equivalent reward within its equivalence class can be expressed using the proposed reparameterization. Introduce the projection operator $f$:
\begin{equation}
    f(r; \pi_\text{ref}, \beta)(x, y) = r(x, y) - \beta\log\sum_{y}\pi_\text{ref}(y\mid x)\exp\left(\frac{1}{\beta}r(x, y)\right)
\end{equation}
This operator effectively adjusts $r$ by subtracting the log-partition function, which depends only on $x$, hence $f(r; \pi_\text{ref}, \beta)(x, y)$ belongs to the same equivalence class as $r(x, y)$. Now, by substituting the reward formulation from Eq.~\ref{eq:main_eq} (valid for all $r$), we see that $f(r; \pi_\text{ref}, \beta)(x, y) = \beta \log \frac{\pi_r(y\mid x)}{\pi_\text{ref}(y\mid x)}$. Therefore, the projection $f$ provides a canonical representative from the equivalence class of $r$ with the desired structure, showing that the reparameterized form suffices without loss of generality.
\end{sproof}

An alternative interpretation of Theorem~\ref{thm:main} is that it identifies the specific reward function, from within each equivalence class, that the DPO reparameterization chooses—namely, the one satisfying the condition:

\begin{equation}\label{eq:lag_p}
     \sum_{y}\underbrace{\pi_\text{ref}(y\mid x)\exp\left(\frac{1}{\beta}r(x, y)\right)}_{=\pi(y\mid x)\text{, by Thm.~\ref{thm:main} reparam.}} = 1,
\end{equation}

which ensures that $\pi(y\mid x)$ defines a legitimate probability distribution (with all probabilities nonnegative and summing to unity). By referencing Eq.~\ref{eq:op_policy}, it becomes evident that Eq.~\ref{eq:lag_p} defines the partition function associated with the optimal policy determined by the reward $r(x, y)$. The principal contribution of the DPO method is the introduction of constraints to the otherwise underdetermined Plackett-Luce (as well as, in particular, the Bradley-Terry) preference models. These constraints maintain the expressivity of the reward models, while ensuring that the optimal policy of Eq.~\ref{eq:op_policy} can be computed in closed form for every prompt $x$.

\subsection{Instability of Actor-Critic Algorithms} Our formalism further enables the identification of instabilities commonly encountered in standard actor-critic methods used in RLHF, like PPO. Consistent with the RLHF workflow, we focus on the policy optimization stage described in Section~\ref{section:prelims}. Parallels can be drawn to the control-as-inference approach~\cite{levine2018reinforcement} for the constrained RL framework described by Eq.~\ref{eq:RL}. Let us consider a parametrized policy $\pi_{\theta}(y\mid x)$, which is optimized by minimizing the divergence $\mathbb{D}_{\text{KL}}[\pi_{\theta}(y|x) \mid \mid \pi^*(y\mid x)]$, where $\pi^*$ is the optimal policy from Eq.~\ref{eq:optimum_model} induced by a learned reward function $r_{\phi}(y, x)$. After some derivation, the corresponding objective becomes:

\begin{equation}\label{eq:AC}
    \max_{\pi_{\theta}}\mathbb{E}_{\pi_{\theta}(y\mid x)}\left[\underbrace{r_{\phi}(x, y) -\beta\log\sum_{y}\pi_\text{ref}(y\mid x)\exp\left(\frac{1}{\beta}r_{\phi}(x, y)\right)}_{f(r_{\phi}, \pi_\text{ref}, \beta)} - \underbrace{\beta\log\frac{\pi_{\theta}(y\mid x)}{\pi_\text{ref}(y\mid x)}}_{\text{KL}}\right]
\end{equation}

The objective under consideration is identical to that of previous studies 
\citep{ziegler2020finetuning, stiennon2022learning, bai2022training, ouyang2022training}, where optimization was performed using a DPO-equivalent reward associated with the reward family $r_{\phi}$. In this framework, the normalization component of $f(r_{\phi}, \pi_\text{ref}, \beta)$ may be interpreted as the soft value function corresponding to the reference policy $\pi_\text{ref}$. Although omitting this normalization term does not alter the optimal solution, its absence can lead to increased variance in the policy gradient estimate, thereby compromising training stability. It is possible to offset the normalization using a learned value function; however, optimizing this auxiliary objective can itself be challenging. Alternatively, earlier approaches have employed reward normalization based on human completion baselines, which act as single-sample Monte Carlo approximations of the normalization term. By contrast, the DPO reparameterization leads to a reward function that operates independently of such baseline estimators.

\section{Experiments}
In this section, we present empirical assessments of DPO’s capacity to train policies from preference signals alone. Initially, in a controlled text-generation scenario, we investigate how effectively DPO balances reward maximization against minimizing the KL-divergence to a reference policy, as compared to established preference optimization techniques such as PPO. Subsequently, we extend our evaluation to encompass larger model architectures and more complex RLHF tasks—including both summarization and dialogue. Remarkably, we observe that with minimal hyperparameter adjustment, DPO often matches or surpasses robust baselines, such as RLHF implemented via PPO, and outperforms approaches that select the highest-reward trajectory out of $N$ samples generated with a learned reward model. Prior to detailing these empirical outcomes, we provide an account of our experimental methodology, with supplementary information furnished in Appendix~\ref{app:exp_details}.

\textbf{Tasks.} We investigate three distinct open-ended text generation tasks in our experiments. Across all tasks, the learning algorithms derive a policy from a collection of preference-labeled examples $\mathcal{D}=\bigl\{x^{(i)}, y_w^{(i)}, y_l^{(i)}\bigr\}_{i=1}^N$. In the case of \textbf{controlled sentiment generation}, $x$ serves as the beginning segment of a movie review drawn from the IMDb dataset \cite{maas-EtAl:2011:ACL-HLT2011}, with the policy required to continue $x$ in a way that exhibits positive sentiment. For this controlled evaluation, we synthetically produce preference pairs by evaluating model outputs with a pre-trained sentiment classifier, ensuring that $p(\text{positive}\mid x,y_w)>p(\text{positive}\mid x,y_l)$. For SFT, GPT-2-large is fine-tuned until convergence using the training split of IMDb reviews (additional details provided in App~\ref{app:sentiment_details}). In the \textbf{summarization} task, $x$ corresponds to a Reddit forum post, and the objective is to generate a summary $y$ capturing the main ideas. Following previous methodologies, we employ the Reddit TL;DR dataset \citep{volske-etal-2017-tl} combined with human preference annotations curated by \citeauthor{stiennon2022learning}. Here, the SFT model is adapted via fine-tuning on human-written forum post summaries\footnote{\url{https://huggingface.co/CarperAI/openai_summarize_tldr_sft}} using the TRLX framework \citep{leandro_von_werra_2023_7790115} for RLHF. The human preference labels are obtained from samples generated by a different SFT model, but one trained similarly, as reported by \citeauthor{stiennon2022learning}. For \textbf{single-turn dialogue}, $x$ is a human prompt—ranging from technical questions in astrophysics to requests for personal guidance. The policy is expected to generate a helpful and engaging response $y$ to this prompt. For this, we rely on the Anthropic Helpful and Harmless dialogue dataset \citep{bai2022training}, comprising 170,000 dialogues between humans and automated systems, each ending with two candidate completions and a human-annotated preference indicating the superior response. As no pre-trained SFT model exists in this context, we create an SFT model by fine-tuning an existing language model using only the preferred response completions.

\textbf{Evaluation.} We adopt two distinct evaluation methodologies in our experimental setup. To systematically assess how well each algorithm optimizes the constrained reward maximization objective, we conduct controlled sentiment generation experiments, measuring algorithm performance via the frontier between the obtained reward and KL-divergence relative to the reference policy; this analysis is feasible because the true reward function—a sentiment classifier—is available. In contrast, in real-world contexts where the reward function is not directly accessible, algorithm performance is gauged by comparing their \textit{win rate} against a baseline policy. For this, GPT-4 serves as a stand-in for human raters, evaluating summary quality and response helpfulness in the summarization and single-turn dialogue tasks, respectively. The reference summary from the test set acts as the baseline for summarization, while the test dataset's preferred response is used as the baseline for dialogue. Previous works suggest that large language models are often more reliable than conventional automated metrics for evaluation \citep{Chen2023ExploringTU}, yet, to substantiate our reliance on GPT-4, we conduct a human subject study detailed in Sec.~\ref{sec:human-judgments}. The results demonstrate that GPT-4's assessments are highly consistent with those made by humans, often achieving parity with or exceeding the level of agreement observed among human annotators.

\textbf{Methods.} Beyond DPO, we benchmark several prevailing methods for aligning language model behavior with human preferences. As baseline strategies, we employ zero-shot prompting with \textbf{GPT-J} \citep{gpt-j} for the summarization task and 2-shot prompting using \textbf{Pythia-2.8B} \citep{biderman2023pythia} in dialogue. We further include the \textbf{SFT} model and a \textbf{Preferred-FT} variant, which is fine-tuned through supervised learning on the selected completion $y_w$—either sourced from the SFT model (for controlled sentiment and summarization tasks) or from a general LM (in single-turn dialogue). Another method we examine is \textbf{Unlikelihood}~\citep{welleck2019neural}, where training maximizes the probability assigned to $y_w$ and simultaneously reduces the likelihood of $y_l$, controlled by an optional coefficient $\alpha\in[0,1]$ applied to the unlikelihood term. We also assess \textbf{PPO} \citep{schulman2017proximal}, which utilizes a reward function learned from preference data, and its oracle variant \textbf{PPO-GT}, which leverages the known reward function available in the controlled sentiment setting. In these experiments, PPO-GT is evaluated using two implementations: an off-the-shelf version \cite{leandro_von_werra_2023_7790115} and a modified version that applies reward normalization and optimized hyperparameters for enhanced effectiveness; these modifications are similarly employed for PPO runs with learned rewards. Lastly, we test the \textbf{Best of $N$} approach, which involves sampling $N$ candidate outputs from the SFT model (or from Preferred-FT for dialogue) and selecting the highest-scoring output according to a reward model derived from preference data. Despite its strong empirical performance, this method dissociates reward model quality from PPO-style optimization and is computationally intensive, as it necessitates generating $N$ completions per query during evaluation.

\subsection{How well can DPO optimize the RLHF objective?}

Typical RLHF approaches employ a reward maximization objective with a KL constraint, which serves to maximize the reward while preventing the policy from straying excessively from a reference policy. Consequently, when evaluating different algorithms, it is essential to jointly consider both the level of reward obtained and the extent of the KL divergence; modest improvements in reward that are accompanied by large increases in KL are often undesirable. Figure~\ref{fig:frontier-tldr-main} illustrates the tradeoff between reward and KL for a range of methods in the sentiment analysis task. For each algorithm, we conduct several independent training runs, each corresponding to a different setting of the conservativeness hyperparameter (target KL values of $\{3,6,9,12\}$ for PPO; $\beta$ values of $\{0.05,0.1,1,5\}$, and $\alpha$ values of $\{0.05,0.1,0.5,1\}$ for unlikelihood; and multiple random seeds for preferred-FT), culminating in a total of 22 experimental runs. Every 100 training iterations, up to convergence, each resulting policy is assessed on a set of held-out prompts by calculating both the mean reward under the actual reward function and the mean sequence-level KL\footnote{Where the sequence-level KL is the sum of KL divergences at each timestep.} with respect to the reference policy, denoted as $\text{KL}\left(\pi\mid \mid \pi_\text{ref}\right)$. Our empirical findings indicate that DPO achieves a significantly more favorable reward-KL frontier than other methods, attaining high rewards with relatively small KL increases. This outcome stands out for several reasons. Notably, although both DPO and PPO seek to optimize the same objective, DPO demonstrates superior efficiency, offering a strictly better tradeoff between reward and KL compared to PPO. Moreover, DPO surpasses PPO’s performance frontier even when PPO is supplied with exact reward information (PPO-GT).

\subsection{Can DPO scale to real preference datasets?}
\label{sec:dpo-real-datasets}
We proceed to assess the effectiveness of DPO fine-tuning using real-world preference data, focusing on tasks such as summarization and single-turn dialogue. In the context of summarization, automatic evaluation metrics like ROUGE often show weak alignment with actual human preferences~\citep{stiennon2022learning}, and earlier studies have demonstrated that leveraging human feedback via PPO leads to more satisfactory summaries. Our comparative analysis involves generating model outputs on the test portion of the TL;DR summarization dataset and determining the average proportion of instances in which model completions outperform references in the test set. We sample responses across a range of temperatures from 0.0 to 1.0, and illustrate the resulting win rates in Figure~\ref{fig:frontier-tldr-main} (right). For a fair comparison, DPO, PPO, and Preferred-FT each fine-tune the identical GPT-J SFT model\footnote{\url{https://huggingface.co/CarperAI/openai_summarize_tldr_sft}}. Our results indicate that at temperature 0.0, DPO attains a win rate of about 61\%, surpassing PPO’s best performance of roughly 57\% at its optimal sampling temperature. DPO also outperforms the best-of-$N$ baseline with respect to peak win rate. It is important to note that DPO's $\beta$ hyperparameter was not extensively optimized in these experiments, suggesting further improvements are possible. Additionally, DPO demonstrates a higher degree of robustness across different sampling temperatures compared to PPO, which exhibits a steep decline in win rate, reverting to the performance of the original GPT-J at higher temperatures. The Preferred-FT approach does not yield substantial improvements over the baseline SFT model. Furthermore, in direct human comparison studies discussed in Section~\ref{sec:human-judgments}, DPO outputs sampled at temperature 0.25 were selected over PPO’s outputs (also at temperature 0) 58\% of the time.

For single-turn dialogue, we assess the various approaches using a subset of the test split from the Anthropic HH dataset \citep{bai2022training} consisting of single human-assistant exchanges. For GPT-4 assessments, we use the dataset’s preferred completions as a reference to calculate the win rate for each approach. Since a standardized SFT model does not exist for this particular setting, we initialize with a pre-trained Pythia-2.8B, fine-tune a reference model with Preferred-FT using the selected completions to ensure distributional alignment, and subsequently apply DPO training. Additionally, we benchmark against the best output from 128 Preferred-FT completions (having established in our experiments that the Best of $N$ baseline stabilizes at $N=128$ for this dataset; refer to Appendix Figure~\ref{fig:best-of-n}) and compare with a 2-shot prompted version of the original Pythia-2.8B. Our findings indicate that DPO matches or surpasses competing methods across optimal sampling temperatures. Furthermore, we evaluate an RLHF model trained with PPO on the Anthropic HH dataset\footnote{\url{https://huggingface.co/reciprocate/ppo_hh_pythia-6B}}, sourced from a reputable implementation\footnote{\url{https://github.com/CarperAI/trlx/tree/main/examples/hh}}, but are unable to identify prompt or temperature settings that outperform the unadapted Pythia-2.8B model. Drawing on both the TL;DR experiment results and the observation that both Best of 128 and PPO target the same reward metric, we treat Best of 128 as a rough indicator of PPO capability. In summary, DPO stands out as the only computationally efficient technique to surpass preferred completions in the Anthropic HH dataset and achieves parity or better relative to the more compute-intensive Best of 128 method. Lastly, Figure~\ref{fig:dialogue-main} demonstrates that DPO reaches peak performance after comparatively few training steps.

\subsection{Generalization to a new input distribution}

To more thoroughly investigate how PPO and DPO handle shifts in data distribution, we evaluate the policies obtained from the Reddit TL;DR summarization task on an out-of-domain dataset—namely, the test split of CNN/DailyMail \citep{nallapati-etal-2016-abstractive}, which comprises news articles. The evaluation utilizes the optimal sampling temperatures determined from TL;DR (0 and 0.25). Table~\ref{tab:ood} displays these findings. We assessed the performance of each method by computing the GPT-4 win rate, comparing system-generated summaries with ground-truth summaries, utilizing the same GPT-4 (C) evaluation prompt as in the Reddit TL;DR task, except that the phrase ``forum post'' was substituted with ``news article''. On this distinct data distribution, DPO significantly surpasses PPO, maintaining its lead. This study offers preliminary support for the notion that DPO policies can generalize to new domains at least as effectively as PPO, despite DPO not utilizing the additional unlabeled Reddit TL;DR prompts available to PPO.

\subsection{Validating GPT-4 judgments with human judgments}
\label{sec:human-judgments}
A human evaluation was conducted to test the reliability of GPT-4's preference judgments, referencing outcomes from the TL;DR summarization experiment and applying two distinct GPT-4 prompt variants. The \textbf{GPT-4 (S)} (simple) prompt solicits an assessment of which summary better captures the key information in a post. The \textbf{GPT-4 (C)} (concise) prompt asks not only about informativeness but also requests a judgment on conciseness; this variant was included because, under the \textbf{GPT-4 (S)} prompt, GPT-4 tends to favor longer and more redundant summaries than those preferred by humans. Complete versions of both prompts are available in Appendix~\ref{app:prompts}. For the evaluation, we selected three system outputs: the highest-performing (DPO, temperature 0.25), the lowest-performing (PPO, temperature 1.0), and an intermediate (SFT, temperature 0.25), ensuring a wide spectrum of summary qualities. All methods were benchmarked against PPO at its best greedy sampling setting. Results indicate that GPT-4's judgments—across both prompt styles—are in alignment with human preferences to a similar degree as human annotators agree amongst themselves, implying that GPT-4 can act as an effective surrogate for human evaluation (though for DPO and PPO-1, due to resource constraints, multiple human ratings were gathered). Notably, the \textbf{GPT-4 (C)} prompt produces win rates most reflective of human assessments; accordingly, this prompt is used for the principal results in Section~\ref{sec:dpo-real-datasets}. Further methodological specifics—such as the rating interface and participant roster—are given in Appendix~\ref{app:human-study}.

\section{Discussion}
Preference-based learning provides a robust and scalable approach for developing advanced, well-aligned language models. In this work, we proposed DPO, a straightforward training methodology that enables learning from preference data without the need for reinforcement learning. Rather than adapting the preference learning task to fit the traditional RL framework and applying standard RL algorithms, DPO establishes a principled correspondence between policy and reward function in language models. This allows direct optimization for human preferences using a basic cross-entropy loss, circumventing reinforcement learning altogether and maintaining generality. DPO requires minimal hyperparameter adjustment and achieves results comparable to, or surpassing, leading RLHF approaches such as those employing PPO. As such, DPO significantly lowers the practical hurdles to leveraging human preferences for training language models.

\textbf{Limitations \& Future Work.} The findings presented here prompt several avenues for further investigation. One open question concerns the out-of-distribution generalization capabilities of the DPO-trained policy in comparison to methods that rely on explicit reward models. Our preliminary analyses indicate that DPO-trained policies can exhibit generalization behavior similar to that observed with PPO-based models, though this warrants more thorough examination. For instance, it remains to be seen if self-labeled data from a DPO policy can facilitate efficient use of unannotated prompts. Furthermore, the occurrence and implications of reward over-optimization within the direct preference optimization framework deserve scrutiny, particularly in light of the modest performance dip shown in Figure~\ref{fig:dialogue-main}-right, which may be a case in point. While our evaluation is limited to models up to 6B parameters, investigating DPO’s scalability to much larger, state-of-the-art models presents an intriguing future direction. On the topic of evaluation methodology, we observe that GPT-4-calculated win rates are sensitive to the input prompt, suggesting a need for further research on strategies for obtaining reliable automated assessments. Beyond human preference learning for language models, DPO offers potential applications in training generative models across a variety of domains.